% Generated draft: requires mathematical and source review.
\documentclass[11pt]{article}
\usepackage[margin=25mm]{geometry}
\usepackage{fontspec}
\setmainfont{Latin Modern Roman}
\usepackage{amsmath,amssymb,mathtools,unicode-math}
\setmathfont{Latin Modern Math}
\usepackage{enumitem,needspace,xurl,hyperref}
\setlength{\parindent}{0pt}
\setlength{\parskip}{5pt}
\setlength{\emergencystretch}{3em}
\newcommand{\N}{\mathbb{N}}
\newcommand{\Z}{\mathbb{Z}}
\newcommand{\Q}{\mathbb{Q}}
\newcommand{\R}{\mathbb{R}}
\newcommand{\C}{\mathbb{C}}
\newcommand{\F}{\mathbb{F}}
\newcommand{\A}{\mathbb{A}}
\newcommand{\PP}{\mathbb{P}}
\newcommand{\E}{\mathbb{E}}
\newcommand{\eps}{\varepsilon}
\newcommand{\dd}{\,\mathrm{d}}
\begin{document}
\section*{UnsolvedMath ID: 30003240}
\textbf{Source problem code:} catalogue entry
\subsection{Definitions and mathematical statement}
The supplied record concerns rings denoted by
\[
H=H_{t_1,t_2}(\C^2),\qquad M=M_{t_1,t_2}(\C^2),
\]
described as rings of entire functions on \(\C^2\) and their meromorphic fraction fields. Here an entire function is a holomorphic function on all of \(\C^2\). For an integral domain \(H\) of such functions, its fraction field is
\[
\operatorname{Frac}(H)=\{f/g:f,g\in H,\ g\ne0\},
\]
with fractions identified by equality after cross multiplication. The record describes \(M\) as this fraction field. It does not define the subscripts \(t_1,t_2\), specify which entire functions belong to \(H\), or explain whether the subscripts also indicate distinguished language symbols.

A relation on a structure is \emph{positive-existentially definable} if it is defined by a first-order formula built from atomic formulas using finite conjunctions, finite disjunctions, and existential quantification. Such a formula uses neither negation nor universal quantification. Definability depends on the chosen language and the permitted parameters.

\textbf{Evaluation-definability questions.} For each of \(A=H\) and \(A=M\), is the evaluation predicate \(\operatorname{Eval}_A\) positive-existentially definable in the corresponding ring language? The original statement specifies that the meromorphic language includes the ordinary valuation predicate. Schematically, if the intended evaluation relation has arity \(r\), the question asks whether there is a positive-existential formula \(\varphi_A(x_1,\ldots,x_r)\) such that
\[
\forall(a_1,\ldots,a_r)\in A^r,\qquad
\operatorname{Eval}_A(a_1,\ldots,a_r)
\Longleftrightarrow
A\models\varphi_A(a_1,\ldots,a_r).
\]
This is a schema, not a complete specification: the record supplies neither the arity nor the interpretation of \(\operatorname{Eval}_A\). It also leaves unspecified the valuation, its predicate, the full language signatures, and whether parameters are allowed. Consequently, no particular evaluation point or specialization map is asserted here.

\textbf{Additional decision questions from the original statement.} Is there an algorithm deciding whether an algebraic equation with coefficients in \(\Z[z]\) has a solution analytic at \(0\), on a neighbourhood of the unit disc, or on all of \(\C\)? These are three separate analytic existence conditions.

For clarity, if an input equation is intended to be
\[
P(z,X_1,\ldots,X_n)=0,
\qquad P\in\Z[z,X_1,\ldots,X_n],
\]
then the usual holomorphic-function interpretation would ask for holomorphic functions \(f_1,\ldots,f_n\) on a suitable open set \(U\subseteq\C\) satisfying
\[
P\bigl(z,f_1(z),\ldots,f_n(z)\bigr)=0
\qquad\text{for every }z\in U.
\]
The three domain conditions would respectively be: some open neighbourhood \(U\) of \(0\); some open neighbourhood \(U\) of the intended unit disc; and \(U=\C\). This explains a possible formal template without fixing the missing input specification. The record does not say whether there is one unknown function or several, whether their number is fixed or part of the input, or whether systems of equations are included. It also does not specify whether the unit disc is open or closed. These choices must be resolved before a precise decision problem can be asserted.

The supplied curation labels the evaluation questions open and states that a negative result in the report concerns a different differential-equation variant. That variant is not specified here, and its negative result is not asserted for the questions above.
\subsection{Short English statement}
The main questions ask whether evaluation can be expressed using positive-existential formulas in the stated entire-function ring and meromorphic fraction field, with a valuation predicate in the meromorphic language. The original record also asks about deciding the existence of analytic solutions of polynomial equations in three domains. Essential definitions and input conventions are missing, so these remain explicitly qualified statement schemas requiring mathematical review.
\subsection{Sources}
\begin{itemize}
\item Supplied local UnsolvedMath registry record, permanent ID \(30003240\), source problem code \texttt{catalogue entry}, title \emph{Definability of Evaluation in Rings of Entire Functions}; collection: Oberwolfach Reports Open Problems. The record gives proposed year 2016 and names no individual proposer.
\item Supplied report citation: \emph{Definability and Decidability Problems in Number Theory} (2017), \emph{Oberwolfach Reports}, pp.\ 2793--2866, DOI \texttt{10.4171/owr/2016/49}. Supplied source link: \url{https://doi.org/10.4171/owr/2016/49}. The background describes this report collection as documenting mathematical workshops and their open-problem discussions.
\item Supplied literature source, labelled ``Oberwolfach Reports 2017'': \url{https://ems.press/content/serial-article-files/46652}.
\item Literature assessment dated \texttt{2026-08-21}: the supplied curation labels the positive-existential evaluation questions open and distinguishes them from a negative result about a different differential-equation variant. This is supplied dated curation, not an independently verified literature assessment. The source background expressly requests independent verification of time-sensitive status and cited claims.
\item The registry supplies status \texttt{open}, statement status \texttt{corrected\_verified}, and a claim of manual semantic, typographic, and LaTeX review. These are source metadata; they do not resolve the missing definitions identified in this draft. The cited report has not been independently consulted for this draft.
\end{itemize}
\end{document}
